% ============================================================
\chapter{Mod\`eles Pr\'e-entra\^in\'es}
\label{ch:pretraines}
% ============================================================

En 2018, trois articles transforment le TAL~: ELMo (Peters et al.), GPT (Radford et al.) et surtout BERT (Devlin et al.). Leur id\'ee commune~: pr\'e-entra\^iner un grand mod\`ele de langage sur des milliards de mots non \'etiquet\'es, puis l'adapter (\emph{fine-tune}) sur la t\^ache cible avec tr\`es peu de donn\'ees supervis\'ees. Ce paradigme \og pr\'e-entra\^inement / adaptation \fg r\'evolutionne la discipline~: les repr\'esentations apprises capturent la syntaxe, la s\'emantique, et m\^eme des connaissances factuelles du monde, rendant les approches sp\'ecialis\'ees ant\'erieures largement obsol\`etes. C'est le d\'ebut de l'\`ere des \emph{foundation models}.

\begin{intuition}[title=\textbf{Le paradigme pr\'e-entra\^inement / adaptation}]
Plut\^ot que d'entra\^iner un mod\`ele de z\'ero pour chaque t\^ache, on
pr\'e-entra\^ine un grand mod\`ele sur des donn\'ees non \'etiquet\'ees, puis on
l'adapte (\emph{fine-tune}) sur des donn\'ees sp\'ecifiques. Ce paradigme,
inaugur\'e par ELMo (2018) et popularis\'e par BERT (2018), a transform\'e le TAL.
\end{intuition}

% ------------------------------------------------------------
\section{Apprentissage de repr\'esentations contextuelles}
\label{sec:pre:contextuelles}
% ------------------------------------------------------------

\begin{definition}[Repr\'esentation contextuelle]
Une \emph{repr\'esentation contextuelle} est une fonction
$\phi : \mathcal{V}^T \times \{1, \ldots, T\} \to \R^d$ qui associe \`a
chaque token $x_t$ un vecteur $\mathbf{h}_t$ d\'ependant du contexte
$(x_1, \ldots, x_T)$ :
\[
\mathbf{h}_t = \phi(\mathbf{x}, t) \neq \phi(\mathbf{x}', t) \quad \text{si } \mathbf{x} \neq \mathbf{x}'
\]
m\^eme si $x_t = x'_t$.
\end{definition}

Cela contraste avec les plongements statiques o\`u $\phi(w)$ est ind\'ependant
du contexte.

% ------------------------------------------------------------
\section{ELMo : contexte par RNN bidirectionnel}
\label{sec:pre:elmo}
% ------------------------------------------------------------

\begin{definition}[ELMo]
ELMo (Embeddings from Language Models, Peters et al., 2018) produit des
repr\'esentations contextuelles en combinant les couches d'un BiLSTM
pr\'e-entra\^in\'e comme mod\`ele de langue bidirectionnel :
\[
\text{ELMo}_t = \gamma \sum_{\ell=0}^{L} s_\ell \mathbf{h}_t^\ell
\]
o\`u $\mathbf{h}_t^\ell$ est l'\'etat cach\'e de la couche $\ell$ \`a la
position $t$, $s_\ell$ sont des poids appris sp\'ecifiques \`a la t\^ache, et
$\gamma$ est un facteur d'\'echelle.
\end{definition}

% ------------------------------------------------------------
\section{BERT : Bidirectional Encoder Representations from Transformers}
\label{sec:pre:bert}
% ------------------------------------------------------------

\subsection{Architecture}

BERT utilise un Transformer encodeur (auto-attention bidirectionnelle) avec
les configurations suivantes :

\begin{center}
\begin{tabular}{lcc}
\toprule
& \textbf{BERT-base} & \textbf{BERT-large} \\
\midrule
Couches $N$ & 12 & 24 \\
Dimension $d$ & 768 & 1024 \\
T\^etes $H$ & 12 & 16 \\
Param\`etres & 110M & 340M \\
\bottomrule
\end{tabular}
\end{center}

\subsection{Objectifs de pr\'e-entra\^inement}

\begin{definition}[Masked Language Model (MLM)]
L'objectif MLM masque al\'eatoirement $15\%$ des tokens d'entr\'ee et
entra\^ine le mod\`ele \`a les pr\'edire :
\[
\mathcal{L}_{\text{MLM}} = -\sum_{t \in \mathcal{M}} \log P(x_t \mid \mathbf{x}_{\setminus \mathcal{M}})
\]
o\`u $\mathcal{M}$ est l'ensemble des positions masqu\'ees. Parmi les tokens
s\'electionn\'es : $80\%$ sont remplac\'es par \texttt{[MASK]}, $10\%$ par un
token al\'eatoire, $10\%$ inchang\'es.
\end{definition}

\begin{definition}[Next Sentence Prediction (NSP)]
L'objectif NSP entra\^ine le mod\`ele \`a pr\'edire si la phrase $B$ suit
la phrase $A$ dans le corpus original :
\[
\mathcal{L}_{\text{NSP}} = -\log P(\text{IsNext} \mid \texttt{[CLS]}, A, \texttt{[SEP]}, B)
\]
\end{definition}

\begin{attention}[title=\textbf{Limites du MLM}]
Le masquage cr\'ee un d\'ecalage entre pr\'e-entra\^inement (tokens
\texttt{[MASK]} pr\'esents) et inf\'erence (pas de \texttt{[MASK]}). Des
travaux ult\'erieurs (RoBERTa, ALBERT) ont montr\'e que le NSP n'est pas
n\'ecessaire et que l'entra\^inement b\'en\'eficie de s\'equences plus longues
et de plus de donn\'ees.
\end{attention}

\subsection{Tokenisation WordPiece}

BERT utilise WordPiece, une variante de BPE qui maximise la vraisemblance :
\[
\text{score}(a, b) = \frac{\text{count}(ab)}{\text{count}(a) \cdot \text{count}(b)}
\]

% ------------------------------------------------------------
\section{GPT : Generative Pre-trained Transformer}
\label{sec:pre:gpt}
% ------------------------------------------------------------

\begin{definition}[GPT]
GPT (Radford et al., 2018) utilise un Transformer d\'ecodeur (auto-attention
causale) pr\'e-entra\^in\'e avec un objectif auto-r\'egressif :
\[
\mathcal{L}_{\text{GPT}} = -\sum_{t=1}^{T} \log P(x_t \mid x_1, \ldots, x_{t-1})
\]
\end{definition}

\begin{center}
\begin{tabular}{lccc}
\toprule
\textbf{Mod\`ele} & \textbf{Param\`etres} & \textbf{Couches} & \textbf{Contexte} \\
\midrule
GPT-1 & 117M & 12 & 512 \\
GPT-2 & 1.5B & 48 & 1024 \\
GPT-3 & 175B & 96 & 2048 \\
GPT-4 & non divulgu\'e & --- & 128K \\
\bottomrule
\end{tabular}
\end{center}

\begin{remarque}
GPT-3 a d\'emontr\'e des capacit\'es \'emergentes en \emph{few-shot learning} :
la capacit\'e de r\'esoudre de nouvelles t\^aches \`a partir de quelques
exemples pr\'esent\'es dans le prompt, sans mise \`a jour des poids.
\end{remarque}

% ------------------------------------------------------------
\section{T5 : Text-to-Text Transfer Transformer}
\label{sec:pre:t5}
% ------------------------------------------------------------

\begin{definition}[T5]
T5 (Raffel et al., 2020) reformule toutes les t\^aches de TAL comme des
probl\`emes texte-vers-texte. Le mod\`ele est un Transformer encodeur-d\'ecodeur
pr\'e-entra\^in\'e avec un objectif de d\'ebruitage (\emph{span corruption}) :
des segments contigus de tokens sont remplac\'es par des sentinelles, et le
d\'ecodeur doit reconstruire les segments masqu\'es.
\end{definition}

\begin{formules}[title=\textbf{Comparaison des architectures}]
\begin{center}
\begin{tabular}{lccc}
\toprule
& \textbf{BERT} & \textbf{GPT} & \textbf{T5} \\
\midrule
Architecture & Encodeur & D\'ecodeur & Enc-D\'ec \\
Attention & Bidirectionnelle & Causale & Bid. + Causale \\
Objectif & MLM + NSP & Auto-r\'egressif & Span corruption \\
Utilisation & Compr\'ehension & G\'en\'eration & Les deux \\
\bottomrule
\end{tabular}
\end{center}
\end{formules}

% ------------------------------------------------------------
\section{Lois d'\'echelle (Scaling Laws)}
\label{sec:pre:scaling}
% ------------------------------------------------------------

\begin{theoreme}[Lois d'\'echelle de Kaplan et al.]
La perte de test $L$ d'un mod\`ele de langue suit des lois de puissance par
rapport au nombre de param\`etres $N$, \`a la taille du dataset $D$ et au
budget de calcul $C$ :
\begin{align}
L(N) &\propto N^{-\alpha_N} & \alpha_N &\approx 0.076 \\
L(D) &\propto D^{-\alpha_D} & \alpha_D &\approx 0.095 \\
L(C) &\propto C^{-\alpha_C} & \alpha_C &\approx 0.050
\end{align}
(pour les mod\`eles de type GPT sur des donn\'ees en anglais).
\end{theoreme}

\begin{remarque}
Les lois Chinchilla (Hoffmann et al., 2022) montrent que l'allocation optimale
du budget calcul suit $N \propto D$ : il faut augmenter la taille du mod\`ele
et des donn\'ees \`a la m\^eme vitesse.
\end{remarque}

% ------------------------------------------------------------
\section{Utilisation avec Hugging Face}
\label{sec:pre:hf}
% ------------------------------------------------------------

\begin{codeblock}[title=\textbf{BERT pour la classification de texte}]
\begin{minted}{python}
from transformers import AutoTokenizer, AutoModelForSequenceClassification
from transformers import pipeline

# Pipeline de sentiment
clf = pipeline("sentiment-analysis", model="nlptown/bert-base-multilingual-uncased-sentiment")
result = clf("Ce cours de NLP est vraiment excellent !")
print(result)
# [{'label': '5 stars', 'score': 0.73}]

# Utilisation directe
tokenizer = AutoTokenizer.from_pretrained("bert-base-uncased")
model = AutoModelForSequenceClassification.from_pretrained("bert-base-uncased", num_labels=2)
inputs = tokenizer("This is a great course!", return_tensors="pt")
outputs = model(**inputs)
logits = outputs.logits  # (1, 2)
\end{minted}
\end{codeblock}

\begin{codeblock}[title=\textbf{GPT-2 pour la g\'en\'eration de texte}]
\begin{minted}{python}
from transformers import pipeline

generator = pipeline("text-generation", model="gpt2")
text = generator(
    "Natural language processing is",
    max_length=50,
    num_return_sequences=1,
    temperature=0.7
)
print(text[0]["generated_text"])
\end{minted}
\end{codeblock}

\begin{codeblock}[title=\textbf{T5 pour le r\'esum\'e}]
\begin{minted}{python}
from transformers import pipeline

summarizer = pipeline("summarization", model="t5-base")
article = """
Natural language processing (NLP) is a subfield of linguistics,
computer science, and artificial intelligence concerned with the
interactions between computers and human language. The goal is to
enable computers to understand, interpret, and generate human language.
"""
summary = summarizer(article, max_length=30, min_length=10)
print(summary[0]["summary_text"])
\end{minted}
\end{codeblock}

% ------------------------------------------------------------
\section{Variantes et \'evolution}
\label{sec:pre:variantes}
% ------------------------------------------------------------

\begin{description}
    \item[RoBERTa] (Liu et al., 2019) : supprime NSP, entra\^ine plus longtemps
          avec plus de donn\'ees et des s\'equences plus longues.
    \item[ALBERT] (Lan et al., 2020) : factorisation des plongements et partage
          de param\`etres inter-couches pour r\'eduire la taille.
    \item[DeBERTa] (He et al., 2021) : attention d\'ecoupl\'ee avec position
          relative.
    \item[LLaMA] (Touvron et al., 2023) : famille de mod\`eles ouverts de
          Meta, avec RMSNorm, RoPE, et SwiGLU.
    \item[Mistral / Mixtral] (2023--2024) : attention group\'ee (GQA),
          m\'elange d'experts (MoE).
\end{description}

% ------------------------------------------------------------
\section{Extraction de features et probing}
\label{sec:pre:probing}
% ------------------------------------------------------------

\begin{definition}[Probing classifiers]
Les \emph{probing classifiers} sont des classifieurs simples (r\'egression
logistique) entra\^in\'es sur les repr\'esentations gel\'ees d'un mod\`ele
pr\'e-entra\^in\'e pour \'evaluer quelles informations linguistiques sont
captur\'ees par chaque couche.
\end{definition}

Des \'etudes ont montr\'e que les couches inf\'erieures capturent la morphologie
et la syntaxe, tandis que les couches sup\'erieures encodent davantage
d'information s\'emantique.

% ------------------------------------------------------------
\section{Exercices}
\label{sec:pre:exercices}
% ------------------------------------------------------------

\begin{exercice}[$\star$]
Comparez le nombre de param\`etres de BERT-base et GPT-2 (117M). Quelles
diff\'erences architecturales expliquent les \'ecarts ?
\end{exercice}

\begin{exercice}[$\star$]
Expliquez pourquoi BERT ne peut pas \^etre utilis\'e directement pour la
g\'en\'eration de texte auto-r\'egressive.
\end{exercice}

\begin{exercice}[$\star\star$]
Impl\'ementez un probing classifier pour \'evaluer si les repr\'esentations
BERT de la couche 6 encodent l'information de POS tagging. Utilisez le
dataset Universal Dependencies.
\end{exercice}

\begin{exercice}[$\star\star$]
Reproduisez l'exp\'erience de masquage de BERT : calculez la perplexit\'e
pseudo-log-vraisemblance $\text{PLL}(\mathbf{x}) = \sum_{t=1}^{T} \log P(x_t \mid \mathbf{x}_{\setminus t})$
sur un ensemble de phrases.
\end{exercice}

\begin{exercice}[$\star\star\star$]
\'Etudiez l'impact des lois d'\'echelle : entra\^inez des mod\`eles de langue
Transformer de tailles croissantes (1M, 10M, 100M param\`etres) sur le m\^eme
corpus et tracez la courbe perte vs.\ param\`etres en \'echelle log-log.
\end{exercice}
